\documentclass[literature.tex]{subfiles}

\begin{document}
\poetry{Kytice}{Karel Jaromír Erben}
\teorieA{
Lyrickoepická skladba básní.
Žánrově jsou každá originál,
tedy příběhy jako \textsc{Zlatý kolovrat} jsou spíše balady,
\textsc{Věšťkyně} je jasná pověst
a \textsc{Záhořovo lože} je legendou.
Celá \textsc{Kytice}, ne~jenom první báseň stejného jména, je dílem poezijním.
Soubor dvanácti děl (plus \textsc{Lilie}, ta~se~udává zvlášť, kvůli svému zvláštnímu stavu).
}{}{
Úsporný jazyk, který zde budí napětí, je také jedním z hlavních rysů balad. Epický příběh je psán spíše kratšími větami. Citoslovce zde nahrazují slovesa. Básně jsou psány gnónickým veršem. \\
\vspace*{1pt}\\
\textbl{Zvukomalba} -- \textsc{Vodník}, \enquote{a na to pole podle skal zelený mužík zatleskal.} \\
\textbl{Epiteton} -- \textsc{Kytice}, \enquote{Ve skrovnou kytici já tě zavážu, ozdobně stužkou ovinu.} \\
\textbl{Anafora} -- \textsc{Zlatý kolovrat}, \enquote{Kde's má Dorničko! Kde jsi? Kde Jsi?
Kde's má roztomilá?}\\
\textbl{Epifora} -- \textsc{Svatební košile}, \enquote{Co to máš na té tkaničce,
na krku na té tkaničce?}\\
\textbl{Epizeuxis} -- \textsc{Holoubek}, \enquote{Běží časy, běží,
Všecko sebou mění.}\\
\textbl{Apostrofa} -- \textsc{Kytice}, \enquote{Mateří-douško vlasti naší milé,
Vy prosté naše pověsti.}\\
\textbl{Metafora} -- \textsc{Vodník}, \enquote{ale kůra srdce tvého
ničím neobměkla!}\\
\textbl{Metonymie} -- \textsc{Vodník}, \enquote{Mladosti mé jarý štěp
přelomil jsi v půli.}\\
\textbl{Personifikace} -- \textsc{Vodník}, \enquote{Bílé šatičky smutek tají,
v perlách se slzy ukrývají.}\\
\textbl{Přirovnání} -- \textsc{Vodník}, \enquote{v sukničku jako z vodních perel,
nechoď, dceruško, k vodě ven.}\\
}{}{}{
Rýmovaný text.
Autor jednotlivé části rozdělil do pozic tak, ži vznikla zrcadlová souměrnost vlastností.
Tedy \textsc{Kytice} a \textsc{Věštkyně} jsou o naději,
\textsc{Polednice} a~\textsc{Vodník} o nadpřirozených bytostech, které působí záporně,
\textsc{Poklad} a \textsc{Dceřina kletba} jsou o narušení vztahu mezi matkou a dítětem,
\textsc{Svatební košile} a \textsc{Vrba} o určité přeměně člověka,
\textsc{Zlatý kolovrat} a~\textsc{Záhořovo lože} mají hned několik společného; vina, pokání, vykoupení;
\textsc{Štědrý den} a \textsc{Holoubek} jsou o~lásce a~smrti,
\textsc{Lilie} se pak někdy dává na místo \textsc{Vrby}, kde také naplňuje příběh o proměně člověka, ale zpravidla se do této dvojice nezapočítává.
Zrovna \textsc{Holoubek} je zároveň přesně v polovině, k čemuž mu dopomáhá bonusová třináctá povídka a představuje pomyslný přechod mezi novými a starými pověstmi.
}{
Pro ukázku jsem vybral \textsc{Záhořovo lože}\\
\vspace*{1pt}\\
Rozlítil se satan a v zlosti své velí: \\
„Vykoupejte jeho v pekelné koupeli!“ \\
Učinila rota dle jeho rozkazu, \\
připravila lázeň z ohně a mrazu: \\
z jedné strany hoří jako uhel vzňatý, \\
z druhé strany mrzne v kámen ledovatý; \\
a když vidí rota míru naplněnu, \\
obrací zmrzlinu opak do plamenů. \\
Strašlivě řve ďábel, jako had se svíjí, \\
a ho pak již smysl i cit pomíjí. \\
Tu pokynul satan, rota odstoupila, \\
a síla zas nová ďábla oživila. \\
Ale když propuštěn opět dýše lehce, \\
krví psané blány přec vydati nechce. -
}
\teorieB{}{
\textbl{Kytice} je příběh o naději.
Matka zemře, děti jí oplakávají, tak se matka přemění na mateřídouškum aby jim mohla nadále být na blízku a mysticky ochraňovat.\par
\textbl{Svatební košile} představuje klasické poučení \enquote{Dej si pozor, co si přeješ, mohlo by se ti to splnit}.
Dívka se modlí k paně Marii za svého milého, který odešel do války a kterého si chtěla vzít.
Milý zaťukal na dveře a vede jí do kostela.
Na cestě jí nutí se zříct Boha a víry, hned několika způsoby.
Ona vždy poslušně splní, jak žádá, s vidinou svatby.
Nakonec jí však dojde, co se děje a před svým nemrtvým snoubencem se schovává v kostele.
Tady je~ještě menší zápas dobra se zlem, kdy se dívka a snoubenec oba snaží přesvědčit mrtvolu v kostele k poslušnosti.
Nemrtvý k oddanosti a panna žádá Boha, aby ho nechal spát.\par
\textbl{Vodník} Dívka i přes matčiny varování jde k rybníku, kde jí chytí vodník a veme si jí za ženu.
Dívka mu porodí dítě, ale stále je nespokojená.
Žádá ho, aby~se~mohla naposledy vidět s maminkou.
Vodník nakonec vyhoví, ale~matka dceru přesvědčí, aby se už nevracela.
Vodník burácí, vyhrožuje, ale~dívka se~nevrací.
Tak jí jako poslední rozloučení nechá přede dveřmi dítě ve dví.\par
\textbl{Záhořovo Lože} je druhou nejdelší básní ve sbírce.
Mladík byl otcem zaprodán ďáblu.
S vírou v Boha se tak dyvává na cestu do pekla.
Na této cestě potkává Záhoře, velkého loupežníka, který zabil všechny pocestné před ním.
Mladík se s ním domluví, že když ho nechá jít, tak mu povypráví, jaké je~to~v~pekle.
Po roce se setkávají a mladík vypráví, jak sám Satan při výhružce křížem rozkázal ďáblu listinu vydat.
Ten odmítl i po mučení.
Až~nakonec Satan rozkázal k něčemu, zvané \enquote{lože Záhořovo}.
Loupežník se zděsí a odhaluje, že to on je Záhoř, pro kterého už tedy v pekle mají přichystané nejhorší mučení ze všech.
Však i ďábel, který vydržel do té doby všechno, se zděsil a úpis vydal.
Mladík mu radí, ať se dá na víru v Boha a prosí o odpuštění.
Také zarazí jeho kyj do země a nařizuje, ať se kaje u toho kyje.
Po letech se tam vrací jako biskup, Záhoř se mezitím proměnil v~pařez.
Oživá při~pohledu na mladíka, který už není mladíkem a tím jejich cesty končí dalším shledáním.
Jejich těla se rozpadnou v prach a jejich duše se změní v bílé holubice směřující k nebi.
}{
Žádný z příběhů nemá přesně určené období ani místo děje.
I když většina se aspoň částečně odehrává na vesnici.
}{}\newpage
\historie{}{}{}{}{}{
Vystudoval filozofii a práva, snažil se o obnovu českého ducha, tak sepisoval české pověsti.
Také usiloval o rovnoprávnost s Němci.
České pověsti sbíral na venkově, kam sám rád jezdil.
}{
Mezi další autorovu tvorbu patří \textsc{Sto prostonárodních pohádek a pověstí slovanských v nářečích původních}, \textsc{Češké pohádky} a \textsc{Vybrané báje a pověsti národní jiných větví slovanských}.
}
%\kritika{}{}{}
\nazor{
Kytice, i když bravurně napsaná a jednotlivé balady jsou velice poutavé, je ještě zajímavější z historického hlediska.
Příběhy jsou o individuální vině a určitém nadpřirozeném, nepřiměřeném trestu, což je velice poutavá zápletka, která funguje.
}
\explain{
\textbl{Legenda} je příběh světce, v českých zemích je význam legendy úzce spjat s katolickými příběhy o svatích mužích a jejich legendárních činnech. Jinak celosvětově je definičně blíže pověsti.
\textbl{Pověst} je příběh, který může být založen na pravdě (vypravěč tomu alespoň věří), má blízko legendě i pohádce (ta ale na pravdě založená nebývá). Díky tomu, že pověsti jsou často o činech sedláků na vesnici je nemožné dostopovat jejich původ a zjistit, zda skutečně reálný základ mají.
\textbl{Balada} je dalším druhem příběhu, který se vyznačuje určitou ponurostí a většinou nešťastným koncem. \\
\textbl{Zvukomalba} je hlásková nápodoba existujícího zvuku, nejčastěji citoslovce, \enquote{kykyryký}.
\textbl{Epiteton} varianty \textit{constans} pro ustálené přirovnání jako \enquote{zelený háj} a \textit{ornans} pro ozdobné přirovnání jako \enquote{zemřelá slova}.
\textbl{Anafora} spočívá v opakování slov, či slovních spojení na začátku, proto \textit{znovuuvedení}.
\textbl{Epifora} spočívá v opakování slov, či slovních spojení na konci, opak anafory.
\textbl{Epizeuxis} spočívá v opakování slov, či slovních spojení hned po sobě, \textit{připojení}.
\textbl{Apostrofa} oslovení jiného diváka, než běžného, k nepřítomné či zemřelé osobě, neživé věci.
\textbl{Metafora} přenášení významu na základě vnější podobnosti, třeba \enquote{zub pily}, či \enquote{kapka štěstí}.
\textbl{Metonymie} přenášení významu na základě jiné, než vnější podobnosti, tedy \enquote{stráž -- střežení}, \enquote{cestující} a \enquote{\textit{Praha} to neschválila}.
\textbl{Personifikace}, \textit{zosobnění}, nejobecněji přenos lidských vlastností na jiné, než lidské bytosti, či humanoidní zevření abstraktních pojmů.
\textbl{Přirovnání} dvou subjektů na základě určité vlastnosti.
\textbl{Gnómický}, nadčasový, aktuální, neměnný, nepodléhající změnám v čase.
}
\wpage
\end{document}